\section*{РЕФЕРАТ}
\addcontentsline{toc}{section}{РЕФЕРАТ}

Объем работы равен \formbytotal{lastpage}{страниц}{е}{ам}{ам}. 
Работа содержит \formbytotal{figurecnt}{иллюстраци}{ю}{и}{й}, 
\formbytotal{tablecnt}{таблиц}{у}{ы}{}, \arabic{bibcount} 
библиографических источников и \formbytotal{числоПлакатов}{лист}{}{а}{ов} 
графического материала.

\textbf{Ключевые слова:} 
мобильное приложение; йога-студия; Android; Kotlin; запись на занятия; 
расписание; профиль пользователя; уведомления; отзывы; статистика; 
графики; SharedPreferences; Material Design.

\textbf{Объект разработки:} 
мобильное приложение «Йога-студия» для платформы Android, 
предназначенное для организации и проведения занятий йогой, 
автоматизации записи клиентов, управления расписанием и 
отслеживания прогресса.

\textbf{Цель работы:} 
разработка программно-информационной системы для организации 
и проведения занятий йогой.

\textbf{Методы и средства:} 
при разработке использовались язык программирования Kotlin, 
средства разработки Android SDK, локальное хранилище SharedPreferences, 
библиотеки MPAndroidChart для построения графиков и CircleImageView 
для отображения аватаров пользователей. Применялись методы 
объектно-ориентированного анализа и проектирования, язык UML 
для моделирования.

\textbf{Результаты:} 
в процессе создания приложения были реализованы модуль авторизации 
и регистрации пользователей, каталог занятий с фильтрацией по городу, 
расписание с календарём, профиль пользователя, система отзывов и 
рейтинга, статистика и графики прогресса, а также уведомления и 
напоминания.

\textbf{Область применения:} 
разработанное приложение может использоваться йога-студиями и 
фитнес-центрами для автоматизации записи клиентов, а также 
отдельными пользователями для отслеживания личного прогресса.